Since fuzzing is a technique that relies heavily on execution speed for each single input, optimizing 
the instrumentation can save time.
In the previously discussed approach there is a major cause of overhead: the point at which the forkserver is executed.
In fact, as can be seen in \cref{fig:dynamicallyLinkedProgramLoad}, since the forkserver is executed by the dynamic linker 
before the execution of the actual target \texttt{\_start} function, for each input the target binary must execute many operations 
before actually reaching the main:
\begin{itemize}
    \item The \_start function must be executed.
    \item The \_\_libc\_start\_main function must be executed, in which the environment is set up for the main execution.
    \item The .init\_array functions are scanned and, if present, executed.
\end{itemize}
The ideal approach is to execute the forkserver at the very beginning of the main function, 
avoiding the execution of all the preamble functions for each input.
This is, in principle, the same concept exposed in \cref{sec:state_of_the_art} when the \textit{deferred forkserver} 
mechanism was showed, with the difference that, in this case, the forkserver entry point is always fixed at the beginning 
of the main.
To achieve this, the relocation mechanism is used again.
In the injected \gls{SO}, the constructor function no longer contains the forkserver, but is used only to 
prepare its execution. The workflow can be seen in \cref{fig:final_forkserver_invocation}.
Precisely, the constructor function will:
\begin{itemize}
    \item set up the \gls{SHM}, reading the \gls{SHM} ID from the environment variable and assigning its address to the exported symbol.\footnote{
        In the next chapter the \gls{SHM} attachment will be discussed in depth. In fact, the \gls{SO} also includes a \textit{standalone mode} 
        in which there is no need for a \gls{SHM}.
    }
    \item assign to a symbol (which will also be exported) the address of the forkserver function that will be executed.
\end{itemize}
Two symbols will be added to the target binary instead of one:
\begin{itemize}
    \item The first is the usual \gls{SHM} address.
    \item The second is a pointer to the forkserver function.
\end{itemize}
These two symbols are then resolved by the dynamic linker at binary startup and their values are written to the section that 
has been added.
In the initial instrumentation, the address of the forkserver function is first dereferenced and then executed.
When the forkserver exits, the rest of the main function is executed, without re-executing all the binary startup functions 
for each input.
\begin{figure}[ht]
    \centering
    \includegraphics[width=400pt]{FinalForkserverExecutionDiagram.png}
    \caption{Optimized forkserver invocation}
    \label{fig:final_forkserver_invocation}
\end{figure}